\chapter{Training Configuration}

\section{Purpose of the Appendix}

\begin{itemize}
    \item Record the exact training configuration used for the reported GRU and LSTM models.
    \item Separate reproducibility details from the narrative explanation in Chapter 3.
    \item Provide a single place to list commands, config files, hyperparameters, and artefact locations.
\end{itemize}

\section{Software Environment}

\begin{itemize}
    \item State the Python version and dependency management tool used for training.
    \item Identify PyTorch, NumPy, SciPy, TensorBoard, and other relevant packages.
    \item Record whether training was run on CPU, GPU, or Apple Silicon acceleration.
    \item Include the exact setup command: \texttt{cd model \&\& uv sync}.
\end{itemize}

\section{Training Command and Configuration Files}

\begin{itemize}
    \item List the final commands used to train or reload each model.
    \item Identify the JSON config files under \texttt{model/Configs/}.
    \item State the output directories under \texttt{model/Results/}.
    \item Note whether the model was trained from scratch or resumed from an existing checkpoint.
\end{itemize}

\section{Hyperparameters}

\begin{itemize}
    \item Report segment length, batch size, maximum epochs, learning rate, optimiser, and weight decay.
    \item Report recurrent warm-up length and update frequency settings.
    \item Report validation frequency, early-stopping patience, and learning-rate scheduler behaviour.
    \item State the loss terms and weights, including ESR and DC loss.
\end{itemize}

\section{Exported Artefacts}

\begin{itemize}
    \item List final \texttt{model.json} and \texttt{model\_best.json} files.
    \item List generated test audio files and \texttt{training\_stats.json}.
    \item Identify which exported JSON files were copied into \texttt{plugin/Resources/Models/}.
    \item Record model metadata: unit type, hidden size, number of layers, skip connection, and sample rate.
\end{itemize}
